\documentclass[11pt,a4paper]{article}
\usepackage{fontspec}
\usepackage{polyglossia}
\setdefaultlanguage{russian}
\setotherlanguage{english}
\setmainfont{PT Serif}[Ligatures=TeX]
\newfontfamily\cyrillicfont{PT Serif}[Ligatures=TeX]
\setmonofont{PT Mono}[Scale=0.85]
\newfontfamily\cyrillicfonttt{PT Mono}[Scale=0.85]
\usepackage[margin=2.2cm]{geometry}
\usepackage{microtype}
\usepackage{booktabs}
\usepackage{tabularx}
\usepackage{hyperref}
\hypersetup{colorlinks=true, linkcolor=blue!50!black, urlcolor=blue!50!black, citecolor=blue!50!black}
\usepackage{xcolor}
\linespread{1.1}
\usepackage{xurl}
\emergencystretch=3em

% ---- числа, полученные в practice.ipynb ----
\newcommand{\Ntotal}{315}
\newcommand{\Ntrain}{150}
\newcommand{\Ntest}{165}
\newcommand{\Nok}{298}
\newcommand{\Nskip}{17}
\newcommand{\Nfeat}{107}


\begin{document}

\section{Материалы и методы}

\subsection{Описание набора данных}

В работе использовался открытый набор данных пациентов с плоскоклеточным раком ротоглотки (oropharyngeal cancer, OPC), подготовленный рабочей группой MICCAI/M.D.~Anderson Cancer Center Head and Neck Quantitative Imaging Working Group для двух соревнований по радиомике, проводившихся на платформе Kaggle in Class в рамках конференции MICCAI~2016: прогнозирования HPV-статуса и прогнозирования локального рецидива первичной опухоли (Local recurrence prediction challenge)~\cite{elhalawani2017}. Набор включает ретроспективные данные пациентов, получавших лечение в MD Anderson Cancer Center (Хьюстон, США) в 2005--2012~гг.; сбор данных был одобрен институциональным этическим комитетом, а в силу ретроспективного характера исследования требование об информированном согласии было снято. Исходные данные в формате DICOM/DICOM-RT и клинические таблицы находятся в открытом доступе в репозитории figshare~\cite{figshare_lrp, figshare_hpv} без регистрации; описание набора опубликовано по лицензии CC~BY~4.0, а файлы метаданных переданы в общественное достояние (CC0). Данные были деидентифицированы авторами набора (DICOM Anonymizer 1.1.6.1 и последующая проверка в ImageJ). В настоящей работе использовалась версия набора для задачи прогнозирования локального рецидива, распространяемая в рамках учебного курса в виде файлов NIfTI и соответствующая соревнованию \url{https://www.kaggle.com/competitions/opc-recurrence}.

Согласно описанию набора~\cite{elhalawani2017}, в исследование включались пациенты с гистологически подтверждённым плоскоклеточным раком ротоглотки (коды МКБ-10 C01, C05.1, C09, C09.0, C10.0, C10.2, C10.3, C10.8), получившие радикальное лечение методом лучевой терапии с модуляцией интенсивности (IMRT) без предшествующего хирургического лечения и с периодом наблюдения более двух лет, с известным HPV/p16-статусом и с доступным диагностическим КТ-исследованием с внутривенным контрастированием, выполненным до начала лечения. Исключались пациенты с не-плоскоклеточным раком и неизвестным HPV-статусом, а также пациенты без пригодного КТ-исследования: без КТ до начала IMRT, с КТ без контрастирования, с артефактами от металлических зубных конструкций или неполным охватом области интереса, а также пациенты, у которых доступное КТ было выполнено уже после индукционной химиотерапии или эксцизионной биопсии. По описанию авторов, из 978 первично отобранных пациентов этим критериям удовлетворяли 315; после исключения 9~случаев возможного дублирования изображений и 18~случаев с нерепрезентативной сегментацией первичной опухоли в опубликованную когорту вошли 288~пациентов. Распространяемая версия набора, однако, содержит данные \Ntotal{} пациентов (\Ntrain{} в обучающей и \Ntest{} в тестовой части соревнования), т.\,е. соответствует этапу отбора до последнего исключения; это расхождение с описанием в статье следует учитывать при сопоставлении результатов. Для пациентов тестовой части значение целевой переменной (локальный рецидив) скрыто организаторами соревнования.

Для всех \Ntotal{} пациентов доступны КТ-исследование и клинические данные; при этом маска первичной опухоли присутствует не у всех пациентов, так как по условиям разметки сегментировались только видимые очаги. В итоговую таблицу были включены \Nok{} пациентов, для которых радиомические признаки первичной опухоли были успешно рассчитаны; \Nskip{} пациентов (10 из обучающей и 7 из тестовой части) были исключены, поскольку для них отсутствует маска первичной опухоли (GTVp), а значит, область интереса для расчёта признаков не определена. Других причин исключения выявлено не было. Таким образом, итоговая выборка включает 140 пациентов обучающей части (из них у 12 зарегистрирован локальный рецидив, что соответствует всем 12 случаям рецидива в исходной обучающей таблице) и 158 пациентов тестовой части. Дополнительно было проверено наличие побитово идентичных КТ-изображений у разных пациентов (по контрольной сумме MD5 массива вокселей) с учётом упомянутых авторами возможных дубликатов; совпадений выявлено не было.

Единственной модальностью медицинской визуализации в наборе является компьютерная томография головы и шеи с внутривенным контрастированием. По данным авторов, 92\,\% исследований выполнены на томографах GE Medical Systems (LightSpeed16 и LightSpeed VCT) при напряжении на трубке 100--154~кВ и матрице $512\times512$; в использованных файлах размер вокселя в аксиальной плоскости составлял 0,39--0,84 (медиана 0,49)~мм, толщина среза "--- 0,39--5,0~мм (медиана 1,0~мм). Изображения ПЭТ и МРТ в набор не входят.

Клинические данные представлены следующими переменными. Идентификатор пациента (\texttt{Patient\_ID}) "--- случайный номер, присвоенный после анонимизации и совпадающий с именем папки с изображениями. Целевая переменная \texttt{Local\_tumor\_recurrence} кодирует наличие (1) или отсутствие (0) локального рецидива, определённого как рецидив опухоли в той же или иной подобласти ротоглотки. Демографические характеристики включают пол (\texttt{Gender}), возраст на момент постановки диагноза в годах (\texttt{Age\_at\_diagnosis}) и расу (\texttt{Race}). Характеристики опухоли описаны латеральностью (\texttt{Tumor\_side}: правая, левая, двусторонняя), подобластью ротоглотки (\texttt{Tumor\_subsite}: корень языка, миндалина, мягкое нёбо, стенка глотки, язычно-глоточная борозда или неуточнённая), категориями T и N и стадией по классификации AJCC/UICC 7-го издания (\texttt{T\_category}, \texttt{N\_category}, \texttt{AJCC\_Stage}), а также степенью дифференцировки опухоли (\texttt{Pathological\_grade}: I--IV, промежуточные I--II и II--III или NA, если оценка невозможна). Анамнез курения представлен статусом курения на момент диагноза (\texttt{Smoking\_status\_at\_diagnosis}: никогда не куривший, бывший, текущий курильщик) и кумулятивной экспозицией в пачко-годах (\texttt{Smoking\_Pack-Years}; одна пачка сигарет в день в течение года). Параметры лечения включают длительность курса лучевой терапии в днях (\texttt{Radiation\_treatment\_course\_duration}), суммарную предписанную дозу в Гр (\texttt{Total\_prescribed\_Radiation\_treatment\_dose}), число фракций (\texttt{\#\_Radiation\_treatment\_fractions}), а также факт проведения индукционной (\texttt{Induction\_Chemotherapy}) и одновременной с облучением (\texttt{Concurrent\_chemotherapy}) химиотерапии. Жизненный статус пациента на момент последнего наблюдения закодирован в \texttt{KM\_Overall\_survival\_censor} (1 "--- жив, 0 "--- умер). Следует отметить, что HPV/p16-статус упоминается в файле описания переменных, но в клинических таблицах данной версии набора отсутствует. В исходных файлах столбец пачко-лет назывался по-разному в обучающей и тестовой таблицах, поэтому при объединении названия были приведены к единому виду; дополнительно был добавлен признак принадлежности пациента к обучающей или тестовой части (\texttt{Fold}).

\subsection{Протокол извлечения радиомических признаков}

Для каждого пациента набор содержит три файла в формате NIfTI (\texttt{.nii.gz}), полученных конвертацией исходных DICOM/DICOM-RT: КТ-изображение (\texttt{CT.nii.gz}) со значениями интенсивности в единицах Хаунсфилда (HU) и две маски сегментации в той же геометрии (размер, размер вокселя, начало координат и ориентация), что и КТ. Маска \texttt{GTVp.nii.gz} является бинарной и описывает макроскопический объём первичной опухоли (gross tumor volume, primary; метка~1), маска \texttt{GTVn.nii.gz} "--- объём метастатически поражённых регионарных лимфоузлов (GTVn), где каждый узел имеет собственную метку (1, 2, \ldots), нумеруемую начиная с наиболее краниального. Маски были построены вручную по аксиальным срезам диагностических КТ с контрастированием двумя радиационными онкологами в программе VelocityAI 3.0.1, без доступа к клиническим данным, в соответствии с рекомендациями ICRU~62/83 (объём определялся как <<макроскопически выявляемая протяжённость и локализация опухоли>>), с последующей проверкой опытным радиационным онкологом~\cite{elhalawani2017}. При разметке учитывались также данные физикального обследования и фиброскопии. У пациентов без поражения лимфоузлов (N0) или после эксцизионной биопсии лимфоузла маска GTVn отсутствует.

Областью интереса для расчёта признаков служила первичная опухоль (GTVp), поскольку именно её локальный рецидив является конечной точкой задачи. Признаки рассчитывались в Python~3.14 с помощью библиотеки PyRadiomics~\cite{vangriethuysen2017} (версия 3.1.1.dev111, сборка из исходного кода репозитория \url{https://github.com/AIM-Harvard/pyradiomics}, коммит \texttt{8ed5793}; опубликованная на PyPI версия 3.1.0 не поддерживает Python~3.14); для чтения изображений использовалась библиотека SimpleITK~2.5.6. Для всех пациентов применялся единый набор параметров, заданный при создании объекта \texttt{RadiomicsFeatureExtractor}. Поскольку толщина среза и размер вокселя различаются между исследованиями, а текстурные признаки чувствительны к пространственному разрешению~\cite{zwanenburg2020}, изображение и маска перед расчётом были приведены к изотропному вокселю $1\times1\times1$~мм (\texttt{resampledPixelSpacing = [1, 1, 1]}): КТ интерполировалось B-сплайном (\texttt{sitkBSpline}), маска "--- методом ближайшего соседа (поведение PyRadiomics по умолчанию); перед ресэмплингом изображение обрезалось вокруг маски (\texttt{preCrop = True}), что не влияет на значения признаков. Дискретизация интенсивностей для текстурных признаков и гистограммных признаков, требующих дискретизации, выполнялась с фиксированной шириной бина 25~HU (\texttt{binWidth = 25}, значение по умолчанию), что для КТ, имеющих абсолютную шкалу интенсивностей, рекомендуется IBSI вместо фиксированного числа бинов~\cite{zwanenburg2020}. Нормализация интенсивностей (\texttt{normalize = False}), удаление выбросов и ресегментация по диапазону HU не применялись; остальные параметры оставлены по умолчанию (метка маски \texttt{label = 1}, расстояние между соседями для матриц смежности \texttt{distances = [1]}, трёхмерный расчёт, симметричная GLCM с усреднением по 13 направлениям, \texttt{minimumROIDimensions = 2}). Признаки рассчитывались только по исходному изображению (тип изображения \texttt{Original}); производные изображения (вейвлет-, LoG-фильтры и др.) не использовались.

Были включены все классы признаков, доступные в PyRadiomics по умолчанию, что дало \Nfeat{} признака на пациента: 14 признаков формы (\texttt{shape}; объём, площадь поверхности, сферичность, главные оси, максимальные диаметры и др.), 18 признаков первого порядка (\texttt{firstorder}; статистики распределения HU внутри опухоли: среднее, медиана, перцентили, энергия, энтропия, асимметрия, эксцесс и др.) и 75 текстурных признаков второго и более высокого порядка: 24 признака матрицы совместной встречаемости уровней серого (GLCM), 16 "--- матрицы длин серий (GLRLM), 16 "--- матрицы размеров зон (GLSZM), 14 "--- матрицы зависимостей уровней серого (GLDM) и 5 "--- матрицы разностей соседних тонов (NGTDM). Определения признаков соответствуют документации PyRadiomics и в большинстве согласуются со стандартом IBSI~\cite{zwanenburg2020}. Признаки GTVn в рамках данного этапа не рассчитывались.

Расчёт выполнялся в цикле по всем \Ntotal{} пациентам; для каждого пациента проверялось наличие КТ и маски GTVp и непустота маски, а любые ошибки на этапе расчёта перехватывались и протоколировались вместе с причиной. К названиям признаков был добавлен префикс \texttt{GTVp\_}. Рассчитанные признаки объединялись с клиническими данными по идентификатору пациента, в результате чего была сформирована сводная таблица из \Nok{} строк (одна строка на пациента) и 127 столбцов: идентификатор, признак части набора, 18 клинических переменных и \Nfeat{} радиомических признаков. Пропущенных или бесконечных значений среди радиомических признаков не было. При контроле качества по плотностным признакам (медиана или среднее HU в GTVp вне диапазона от $-100$ до $200$~HU) и последующем визуальном осмотре выявлены три пациента тестовой части с вероятными погрешностями маски: у пациентов 219 и 289 в GTVp попадает просвет дыхательных путей (10-й перцентиль около $-960$~HU; медиана у пациента 289 составляет $-624$~HU), у пациента 257 контур GTVp на части срезов располагается на костных структурах (медиана $+223$~HU). Эти пациенты сохранены в таблице без изменения масок, поскольку задачей этапа было воспроизводимое извлечение признаков по исходной разметке, однако они требуют внимания на этапе моделирования. Кроме того, признаки формы, вычисляемые через собственные значения ковариационной матрицы координат вокселей (\texttt{MajorAxisLength}, \texttt{MinorAxisLength}, \texttt{LeastAxisLength}, \texttt{Elongation}, \texttt{Flatness}), в используемой сборке PyRadiomics возвращались в виде комплексных чисел с нулевой мнимой частью; в таблицу записывалась их действительная часть.

\begin{thebibliography}{9}
\bibitem{elhalawani2017}
Elhalawani H., Mohamed A.\,S.\,R., White A.\,L. et al. Matched computed tomography segmentation and demographic data for oropharyngeal cancer radiomics challenges // Scientific Data. 2017. Vol.~4. Art.~170077. \href{https://doi.org/10.1038/sdata.2017.77}{doi:10.1038/sdata.2017.77}.

\bibitem{figshare_lrp}
Fuller C., Mohamed A., Elhalawani H. Determine from CT data whether a tumor will be controlled by definitive radiation therapy [collection] // figshare. 2017. \href{https://doi.org/10.6084/m9.figshare.c.3757385.v1}{doi:10.6084/m9.figshare.c.3757385.v1}.

\bibitem{figshare_hpv}
Fuller C., Mohamed A., Elhalawani H. Predict from CT data the HPV phenotype of oropharynx tumors; compared to ground-truth results previously obtained by p16 or HPV testing [collection] // figshare. 2017. \href{https://doi.org/10.6084/m9.figshare.c.3757403.v1}{doi:10.6084/m9.figshare.c.3757403.v1}.

\bibitem{vangriethuysen2017}
van Griethuysen J.\,J.\,M., Fedorov A., Parmar C. et al. Computational radiomics system to decode the radiographic phenotype // Cancer Research. 2017. Vol.~77, no.~21. P.~e104--e107. \href{https://doi.org/10.1158/0008-5472.CAN-17-0339}{doi:10.1158/0008-5472.CAN-17-0339}. Документация: \url{https://pyradiomics.readthedocs.io}.

\bibitem{zwanenburg2020}
Zwanenburg A., Vallières M., Abdalah M.\,A. et al. The Image Biomarker Standardization Initiative: standardized quantitative radiomics for high-throughput image-based phenotyping // Radiology. 2020. Vol.~295, no.~2. P.~328--338. \href{https://doi.org/10.1148/radiol.2020191145}{doi:10.1148/radiol.2020191145}.
\end{thebibliography}

\end{document}
